\section{Related Work}
\label{sec:related_work}

\subsection{UAV Small-Object Evidence Generation}
\label{sec:rw_uav_small_object}

UAV small-object detection has increasingly focused on preserving shallow
spatial detail while maintaining enough semantic context for classification.
Several recent YOLO-family methods add high-resolution branches, lightweight
feature enhancement, attention, or frequency-sensitive modules to improve tiny
target recall under aerial viewpoints
\cite{hu2024universalyolo,jiang2024mffsodnet,li2024dsyolov5s,bai2025sffefyolo,luo2026fmfnyolo,xie2026dgtsod,chao2025bpdyolo,han2025lrds}.
Other detectors study deformable attention, transformer-style assignment, or
position-sensitive refinement for dense UAV scenes
\cite{hu2025csfprrtdetr,hu2026hfdifine,li2026dfformer,chen2026ufodetr,tan2025rtuavsod,yang2026uavdet}.
Complementary UAV studies use dual-domain attention, coarse-fine feature
alignment, cross-resolution semantic alignment, context modeling, rotated
boxes, segmentation, multimodal fusion, hardware-aware acceleration, and
benchmark comparisons
\cite{wang2026d2adetector,bi2025cfia,tu2026dcocrsa,han2022csadet,ye2022glfnet,lin2025fgunet,tang2026affnet,ghanta2026shield8uav,sadeghibakhi2025comparison}.
These works motivate two design choices in SAFR-YOLO: high-resolution P2/P3/P4
prediction paths for tiny objects and lightweight refinement that preserves
local detail without replacing the YOLO backbone with a large transformer.

Our goal differs from only improving the detector score. We treat detection as
the first evidence-generation stage in a larger multi-UAV system. Therefore the
detector must output not only boxes and classes, but also confidence,
uncertainty, view identity, pose/time metadata, and descriptors that can be
associated across views. This EvidenceToken interface is the bridge from 2D
small-object detection to 3D evidence completion.

\subsection{Multi-View 3D Completion For Aerial Scenes}
\label{sec:rw_3d_completion}

Multi-view recognition and UAV search systems show that heterogeneous
viewpoints must be aligned rather than averaged blindly
\cite{zhang2024attention4align,cai2025neusis,maletic2025spatialsemantic}.
Active UAV reasoning also has physical costs: flight time, energy, bandwidth,
freshness, and safety \cite{jeon2025symbolicrlbdi,emami2026frsicl}. Neural
scene representations and 3D Gaussian-style renderers can reconstruct or render
novel views from sparse image collections, but UAV small-object perception
needs object-centered evidence under viewpoint ambiguity. In our setting, 3D is
not used as a standalone visual effect; it is used to ask whether additional
viewpoints, restored evidence, or approximate geometry can reduce class
ambiguity for small objects.

The MarineCity component of CoM3D-ACE is therefore framed as an evidence
completion protocol. It consumes real-Cesium RGB/depth/pose captures and
detector EvidenceTokens, builds cross-view object hypotheses, and verifies that
Nerfacto, Instant-NGP, and 3DGS-style runners can produce held-out validation
metrics on the exported capture package. We report these rows as system-level
runner-family evidence rather than as a final optimized 3D benchmark.

\subsection{Graph-Grounded Reasoning And Re-Observation}
\label{sec:rw_reasoning}

Large language and vision-language models can help explain ambiguous visual
evidence and support UAV planning, communication, and visual grounding
\cite{tian2025uavsmeetllms,emami2026llmuavsurvey,chen2025uavvisualgrounding}.
However, an unconstrained model should not directly control a UAV mission.
CoM3D-ACE uses a restricted graph-grounded interface: the detector and 3D stages
produce structured evidence, the reasoner receives only the high-ambiguity graph
subproblem, and a verifier checks graph consistency, geometric plausibility, and
action validity. This design lets the reasoner suggest actions such as
finalize, reject, monitor, or targeted re-observation while keeping the final
decision tied to explicit evidence. We call this module the
\textbf{AeroGraph Reasoner}.

Developer studies, safe-control frameworks, and prompt-security analyses also
motivate constrained outputs and safety verification
\cite{chen2025llmuavprojects,zhao2026sagellm,emami2026promptstoprotection}.
The current paper reports AeroGraph as a system interface and verifier-controlled
reasoning module, while keeping language-model response quality separate from
the detector and 3D system-validation claims.
